\documentclass[11pt]{article}
\usepackage{amsmath,amssymb,amsthm,amsfonts}
\usepackage{hyperref}
\usepackage{geometry}
\geometry{margin=1in}

\title{Phase Mirror as Constitutional Geometry:\\
Necklace Topologies, Gurau Degree, and Branched-Polymer Avoidance in Governance Design}
\author{Citizen Gardens \\ The Foundation of Multiplicity}
\date{April 28, 2026}

\begin{document}
\maketitle

\begin{abstract}
This defensive publication formalizes a universal cognitive architecture---Phase Mirror---as a lattice-theoretic and topological framework for governance, and shows how it is literally continuous with tensorial quantum gravity constructions via the Euler characteristic and the Gurau degree. The report develops a small family of toy models (2--4 decision ``necklaces'') that demonstrate how adding decisions with sufficient closed accountability loops generates manifold-like governance topologies, whereas adding decisions without such loops yields branched-polymer pathologies. The publication also specifies a practical ``Phase Mirror spec'' for safe scaling of high-stakes decisions in socio-technical systems.
\end{abstract}
\newpage
\tableofcontents

\section{Executive Summary}

\subsection{Core Claim}

Phase Mirror is presented as a universal cognitive perspective and invariant-consistency oracle for socio-technical controls: a structural diagnostic that applies to any system where decisions are made under abstraction rather than binding. In this report we:

\begin{itemize}
  \item Formalize Phase Mirror governance lattices in terms of vertices, edges, and faces, and link their topology to the Euler characteristic \(\chi\).
  \item Show how these lattices map directly onto rank-\(d\) colored tensor graphs equipped with the Gurau degree \(G\), the central invariant in certain Group Field Theory (GFT) models of quantum gravity.
  \item Construct explicit toy models (2-, 3-, and 4-vertex ``necklace'' graphs) that encode high-stakes enterprise decisions and compute their Gurau degree to diagnose operational viability.
  \item Contrast these necklace configurations with branched-polymer (tree-like) structures, and characterize the transition between manifold-like and pathological governance.
  \item Specify an actionable Phase Mirror deployment spec: how to add new decisions while maintaining manifold-like topology and avoiding branched-polymer enterprise architecture.
\end{itemize}

The central novelty is the tight, operational identification between:
\begin{enumerate}
  \item Governance topology (Phase Mirror lattice of agents, policies, constraints, contracts, SLAs, and kill-switches),
  \item Algebraic-topological invariants (\(\chi\), Betti numbers, homology),
  \item Tensorial quantum gravity invariants (Gurau degree, amplitude convergence, branched polymers vs.\ necklace interactions).
\end{enumerate}

\subsection{Defensive Publication Scope}

This document is intended as prior art that:
\begin{itemize}
  \item Fixes the correspondence between enterprise governance lattices and tensorial quantum gravity graphs.
  \item Encodes specific recipes for constructing necklace vs.\ branched-polymer governance topologies.
  \item Documents a concrete design pattern for adding decisions (vertices) together with accountability loops (faces) such that Gurau degree remains in a safe (non-divergent) regime.
  \item Provides an implementation sketch (in pseudocode / Python-style code) for computing Gurau degree from a governance lattice description.
\end{itemize}

\section{Foundational Geometry: Euler Characteristic and Governance}

\subsection{Euler Characteristic as Structural Invariant}

The Euler characteristic \(\chi\) of a 2-dimensional CW-complex is
\[
\chi = V - E + F,
\]
where \(V\) is the number of vertices, \(E\) the number of edges, and \(F\) the number of faces. For polyhedral surfaces, \(\chi\) is a topological invariant: it remains constant under continuous deformations (bending, stretching, wobbling) that do not tear or glue the surface.

For compact orientable surfaces of genus \(g\), the Euler characteristic satisfies
\[
\chi = 2 - 2g.
\]
This gives:
\begin{itemize}
  \item Sphere (\(g = 0\)): \(\chi = 2\).
  \item Torus (\(g = 1\)): \(\chi = 0\).
  \item Double torus (\(g = 2\)): \(\chi = -2\).
\end{itemize}

Deviation of \(\chi\) from an expected value constitutes mathematical evidence of structural holes. For example, if a shape is claimed to be sphere-like but has \(\chi = 0\), it must possess at least one handle/hole (toroidal behavior).

\subsection{Phase Mirror Lattice Interpretation}

In the Phase Mirror architecture, the Euler formula is reinterpreted as a design blueprint for governance structures:
\begin{itemize}
  \item Vertices (\(V\)) become the autonomous agents, corporate policies, and risk constraints.
  \item Edges (\(E\)) become explicit bindings: technical specifications, liability contracts, SLAs, and kill-switches.
  \item Faces (\(F\)) become closed loops of accountability and verification---full cycles where requirement, implementation, runtime behavior, and oversight are structurally connected.
\end{itemize}

The arithmetic \(\chi = V - E + F\) encodes how connectivity and enclosure co-determine systemic properties. Missing edges or faces manifest as holes in the governance manifold:
\begin{itemize}
  \item \(\chi\) lower than the design target indicates missing closure (faces) or surplus bindings (edges) that create holes and non-contractible loops of behavior.
  \item \(\chi\) higher than expected indicates spurious closure or under-connected components.
\end{itemize}

\section{Quantum Gravity Correspondence: Gurau Degree}

\subsection{Colored Tensor Graphs and Gurau Degree}

In Group Field Theory and related tensor models of quantum gravity, spacetime is represented by colored tensor graphs. For a rank-\(d\) model, one considers a graph with:
\begin{itemize}
  \item Vertices \(V\) (quantum geometric events or ``atoms of spacetime'').
  \item Internal faces \(F_{\text{int}}\) (closed strands corresponding to internal connectivity).
\end{itemize}

The Gurau degree \(G\) of such a graph, in the specialized form relevant here, can be taken as
\[
G = d\,(V - 1) - (d - 1)\,F_{\text{int}},
\]
where \(d\) is the dimension parameter (for spacetime, often \(d = 4\)).

Qualitative interpretation:
\begin{itemize}
  \item Negative or zero \(G\): amplitudes are well-controlled; the configuration supports a viable continuum limit (manifold-like behavior).
  \item Positive \(G\): amplitudes diverge; the configuration tends toward branched-polymer pathology, failing to support smooth geometry under scaling.
\end{itemize}

\subsection{Governance Lattices as Tensor Graphs}

The correspondence fixes:
\begin{itemize}
  \item Governance vertices \(V\): decision events, commitment points, or policy/risk constraint atoms where authoritative determination occurs.
  \item Internal lines/edges: connectivity that makes these events part of a coherent governance whole (authority delegations, information flows, binding contracts).
  \item Internal faces \(F_{\text{int}}\): closed loops of accountability and verification; cycles where decision, implementation, monitoring, and enforcement complete a loop.
\end{itemize}

The remarkable claim is that \emph{when Phase Mirror measures governance vertices and edges, it is literally computing the Gurau degree of the enterprise architecture viewed as a tensorial graph.} The sign and magnitude of \(G\) then become measures of operational viability under stress, scaling, and perturbation.

\section{Toy Models: From Two to Four Vertices}

\subsection{Two-Vertex Necklace}

\subsubsection{Quantum Graph}

Take a minimal nontrivial rank-\(4\) colored tensor graph:
\begin{itemize}
  \item \(d = 4\).
  \item \(V = 2\) vertices.
  \item \(F_{\text{int}} = 4\) internal faces (a simple ``necklace'': multiple strands forming closed loops between the two vertices).
\end{itemize}

Compute:
\[
G_{2} = 4(2 - 1) - 3 \cdot 4 = 4 - 12 = -8.
\]

This strongly negative Gurau degree indicates a well-controlled, manifold-like configuration.

\subsubsection{Governance Embedding}

Interpret the two vertices as:
\begin{itemize}
  \item \(V_1\): Go/No-Go for launching an agentic AI product in healthcare.
  \item \(V_2\): Scale-up authorization for national deployment after pilot.
\end{itemize}

Define four accountability loops (internal faces) that enclose both decisions:
\begin{enumerate}
  \item Regulatory–Legal loop.
  \item Technical–Controls loop.
  \item Risk–Capital loop.
  \item Ethics–Patient loop.
\end{enumerate}

Each loop is a closed cycle of vertices and edges: policies, agents, technical configs, runtime evidence, and enforceable kill-switches. This yields rich accountability closure for only two decision vertices, matching the negative \(G\).

\subsection{Deleting One Loop: Emergence of a Hole}

Remove the Risk–Capital loop:
\begin{itemize}
  \item Vertices remain \(V = 2\).
  \item Internal faces drop to \(F_{\text{int}} = 3\).
\end{itemize}

The Gurau degree becomes:
\[
G_{\text{cut}} = 4(2 - 1) - 3\cdot 3 = 4 - 9 = -5.
\]

We observe:
\begin{itemize}
  \item \(G\) is still negative, but less so; amplitude control weakens.
  \item A specific structural hole appears: decisions can now create unpriced financial risk that is not structurally looped back into risk appetite, capital buffers, or insurance.
\end{itemize}

In Euler terms, with the same \(V\) but fewer faces \(F\), \(\chi\) shifts downward, consistent with a missing face and the creation of a hole in the governance manifold along the financial axis.

\subsection{Three-Vertex Necklace}

Add a third decision:
\begin{itemize}
  \item \(V_3\): Cross-border expansion into a new jurisdiction.
\end{itemize}

Consider a 3-vertex necklace:
\begin{itemize}
  \item \(d = 4\).
  \item \(V = 3\).
  \item \(F_{\text{int}} = 6\), e.g.\ multiple loops connecting \((V_1, V_2)\) and \((V_2, V_3)\).
\end{itemize}

Then:
\[
G_{3} = 4(3 - 1) - 3\cdot 6 = 8 - 18 = -10.
\]

Adding one vertex and two faces makes \(G\) more negative, meaning that the system remains well-behaved under scaling so long as new decisions are accompanied by new loops.

\subsection{Four-Vertex Necklace}

Add a fourth decision:
\begin{itemize}
  \item \(V_4\): Expansion into an adjacent vertical (e.g.\ finance) with the same AI stack.
\end{itemize}

Construct a 4-vertex necklace:
\begin{itemize}
  \item \(d = 4\).
  \item \(V = 4\).
  \item \(F_{\text{int}} = 9\) (multiple local loops plus at least one global loop enclosing all four decisions).
\end{itemize}

Then:
\[
G_{4} = 4(4 - 1) - 3\cdot 9 = 12 - 27 = -15.
\]

As we scale 2 \(\to\) 3 \(\to\) 4 vertices with rich loops, \(G\) becomes more negative, consistently manifold-like. The governance manifold thickens instead of collapsing into a tree.

\subsection{Four-Vertex Branched Polymer}

Contrast with a 4-vertex tree:
\begin{itemize}
  \item \(V = 4\): D1, D2, D3, D4.
  \item \(F_{\text{int}} \approx 0\): no genuine closed loops, only a chain of authority.
\end{itemize}

Then:
\[
G_{\text{tree}} = 4(4 - 1) - 3\cdot 0 = 12.
\]

Here \(G\) is positive; the configuration is in the branched-polymer regime:
\begin{itemize}
  \item Topology is fragile and tree-like.
  \item There are no nontrivial cycles, only paths.
  \item Under scaling, amplitudes diverge: the governance system fails under stress or becomes ungovernable.
\end{itemize}

\section{Phase Mirror Spec: Adding Decisions Without Branched-Polymer Drift}

This section gives a concise specification for how to add high-stakes decisions while preserving manifold-like topology.

\subsection{Invariant Conditions}

Given a set of decision vertices \(V\) in a program:
\begin{itemize}
  \item Maintain a minimum internal-face density: for some \(k \ge 3\), aim for
  \[
  F_{\text{int}} \ge k\,V.
  \]
  \item Monitor Gurau degree:
  \[
  G = d(V - 1) - (d - 1)F_{\text{int}},
  \]
  and ensure \(G\) remains negative or zero and becomes more negative as \(V\) grows.
  \item Maintain at least one global loop that encloses all high-stakes decisions in the program (e.g.\ a global ethics or regulatory loop).
\end{itemize}

\subsection{Operational Rules}

Every time a new decision vertex \(D_n\) is introduced:

\begin{enumerate}
  \item \textbf{Classify change:} Determine which dimensions are impacted:
  \begin{itemize}
    \item Regulatory/legal regime.
    \item Technical/control complexity.
    \item Ethical surface area.
    \item Risk/capital exposure.
  \end{itemize}

  \item \textbf{Add loops for each impacted dimension:} For each impacted dimension, add or upgrade at least one closed loop (face) that:
  \begin{itemize}
    \item Contains \(D_n\) and at least one earlier decision.
    \item Is built from explicit edges: specs, contracts, SLAs, kill-switches.
  \end{itemize}

  \item \textbf{Verify non-vibe status:} For each loop, ensure:
  \begin{itemize}
    \item A formal specification exists (what is done and how it is traced).
    \item A quorum is defined (who must approve or act).
    \item There is a lever with an owner, a metric/threshold, and a time horizon.
  \end{itemize}

  \item \textbf{Runtime activation check:} Require historical or simulated evidence that the loop can be activated by runtime traces (telemetry, incidents) and produce concrete structural changes (slowed rollout, rollback, scope reduction, shutdown).

  \item \textbf{Disallow pure hierarchy as ``governance'':} If \(D_n\) increases stakes but only adds more senior sign-off in a chain, with no new loops, flag this as a branched-polymer move requiring remediation.
\end{enumerate}

\subsection{Documented vs.\ Runtime Topology}

Maintain a dual view:
\begin{itemize}
  \item \emph{Documented lattice}: topology inferred from org charts, policies, contracts, and diagrams.
  \item \emph{Runtime lattice}: topology inferred from actual flows of decisions, alerts, incidents, and interventions.
\end{itemize}

Then:
\begin{itemize}
  \item Compute \(\chi_{\text{doc}}, G_{\text{doc}}\) on the documented lattice.
  \item Compute \(\chi_{\text{run}}, G_{\text{run}}\) on the runtime lattice.
  \item If \(\chi_{\text{run}} < \chi_{\text{doc}}\) or \(G_{\text{run}} > G_{\text{doc}}\), the system has paper loops (faces that are nominally defined but never operationally close). These must be treated as structural holes until runtime closure is demonstrated.
\end{itemize}

\section{Mathematical Overview}

\subsection{Homology and Holes}

Let \(X\) be a finite CW-complex representing a governance lattice. Its homology groups \(H_k(X)\) capture \(k\)-dimensional holes:
\[
\chi(X) = \sum_{k} (-1)^k \beta_k, \quad \beta_k = \text{rank}\,H_k(X).
\]

In governance:
\begin{itemize}
  \item \(\beta_0\): number of connected components (disconnected silos).
  \item \(\beta_1\): 1-dimensional loops not contractible to a point (non-contractible oversight loops or contradictory accountability structures).
  \item \(\beta_2\): 2-dimensional cavities (coverage voids where oversight domains fail to intersect).
\end{itemize}

Deviation of \(\chi\) from target---together with elevated \(\beta_1\) or \(\beta_2\)---identifies structural holes and contradictions.

\subsection{Gurau Degree and Renormalization}

In GFT, the Gurau degree \(G\) governs amplitude behavior under renormalization. As governance is scaled (more decisions, broader scope), an analogous renormalization can be conceptualized:
\begin{itemize}
  \item Coarse-grain sequences of decisions into effective vertices.
  \item Track how \(G\) changes under this coarse-graining.
  \item Systems where \(G\) flows toward positive values under such transformations are governance analogs of non-renormalizable theories: they become more pathological as we zoom out.
\end{itemize}

Conversely, negative or stable \(G\) under coarse-graining indicates governance architectures that maintain their structural integrity across scales.

\section{Implementation Sketch}

Below is a minimal Python-style snippet showing how one might compute a Gurau-like degree for a governance lattice encoded as a simple graph structure. This is illustrative, not fully general.

\begin{verbatim}
from dataclasses import dataclass, field
from typing import List, Set, Tuple

@dataclass
class GovernanceGraph:
    # Each decision is a vertex
    decisions: Set[str] = field(default_factory=set)
    # Edges represent bindings (specs, contracts, SLAs, kill-switches)
    edges: Set[Tuple[str, str]] = field(default_factory=set)
    # Faces are closed loops of vertices (e.g., cycles of accountability)
    faces: List[List[str]] = field(default_factory=list)

    def add_decision(self, name: str):
        self.decisions.add(name)

    def add_edge(self, a: str, b: str):
        self.edges.add(tuple(sorted((a, b))))

    def add_face(self, cycle: List[str]):
        # Assume cycle is a closed loop of decision names
        if len(cycle) > 2 and cycle[0] == cycle[-1]:
            self.faces.append(cycle)

    def gurau_degree(self, d: int = 4) -> int:
        V = len(self.decisions)
        # Internal faces here are approximated by faces list length
        F_int = len(self.faces)
        G = d * (V - 1) - (d - 1) * F_int
        return G

    def euler_characteristic(self) -> int:
        V = len(self.decisions)
        E = len(self.edges)
        # Faces counted as unique cycles; this is simplified
        F = len(self.faces)
        return V - E + F

# Example: 4-vertex necklace
g = GovernanceGraph()
for d in ["D1", "D2", "D3", "D4"]:
    g.add_decision(d)

# Add edges in a chain for simplicity; in reality you'd add all bindings
g.add_edge("D1", "D2")
g.add_edge("D2", "D3")
g.add_edge("D3", "D4")

# Add multiple faces to approximate loops
g.add_face(["D1", "D2", "D1"])   # Reg-Legal for D1-D2
g.add_face(["D2", "D3", "D2"])   # Reg-Legal for D2-D3
g.add_face(["D3", "D4", "D3"])   # Reg-Legal for D3-D4
g.add_face(["D1", "D2", "D1"])   # Technical-Controls D1-D2
g.add_face(["D2", "D3", "D2"])   # Technical-Controls D2-D3
g.add_face(["D1", "D4", "D1"])   # Global Ethics loop D1-D4

G = g.gurau_degree(d=4)
chi = g.euler_characteristic()
print("Gurau degree:", G)
print("Euler characteristic:", chi)
\end{verbatim}

A production implementation would:
\begin{itemize}
  \item Enforce proper cycle detection for faces.
  \item Distinguish internal vs.\ external faces.
  \item Support multi-layer graphs (requirements, configs, code, runtime).
  \item Compute homology groups via a dedicated algebraic topology library.
\end{itemize}

\section{Recommendations and Extensions}

\subsection{Fastest Path to Validation}

To operationalize and validate this framework:

\begin{enumerate}
  \item \textbf{Select a high-stakes program:} e.g., deployment of an agentic AI in a regulated sector.
  \item \textbf{Encode the governance lattice:} Identify decision vertices, explicit bindings, and closed accountability loops.
  \item \textbf{Compute empirical \(\chi\) and \(G\):} Using a tooling stack that represents the lattice and derives approximate Euler characteristic and Gurau degree.
  \item \textbf{Run delta experiments:} 
  \begin{itemize}
    \item Remove or ``vibe-ify'' loops and measure predicted shifts in \(\chi\) and \(G\).
    \item Compare predictions with observed performance under stress tests and scaling.
  \end{itemize}
  \item \textbf{Iterate into design patterns:} Distill reusable patterns: ``minimal necklace for a new jurisdiction,'' ``safe scale-up from regional to national,'' etc.
\end{enumerate}

\subsection{Further Work}

Potential extensions include:
\begin{itemize}
  \item Multi-scale models using persistent homology to track how governance holes appear or disappear across decision timescales.
  \item Direct mapping from governance tensor graphs to numerical GFT simulations, to leverage existing renormalization group tools in quantum gravity.
  \item Integration with causal modeling: representing causal graphs as lattice substructures whose topology contributes to the overall \(G\) and \(\chi\).
\end{itemize}

\section{Conclusion}

This defensive publication captures an integrated framework where:
\begin{itemize}
  \item Governance structures are treated as topological objects with a constitutional geometry.
  \item The Euler characteristic and Gurau degree provide rigorous diagnostics for structural holes and scalability pathologies.
  \item Necklace vs.\ branched-polymer patterns distinguish resilient manifold-like organizations from fragile tree-like ones.
  \item A concrete spec is given for adding decisions while preserving a safe topological regime.
\end{itemize}

By anchoring governance design in invariants that already support the physics of spacetime, this framework aims to ground socio-technical control in a robust, mathematically stable geometry.

\appendix
\section*{Mathematical Appendix}
\addcontentsline{toc}{section}{Mathematical Appendix}

This appendix consolidates the explicit definitions, proofs, and operator-norm bounds that underwrite the governance–topology correspondence and the Gurau-degree diagnostics used in the main text. The goal is to make the “governance as geometry” claims fully explicit in standard mathematical language.

\section{Topological Model of Governance Lattices}

\subsection{Governance as a Finite CW-Complex}

We model a governance architecture as a finite 2-dimensional CW-complex (or more generally a finite CW-complex truncated at dimension 2 for the purposes of this work).

\begin{definition}[Governance lattice]
A \emph{governance lattice} is a finite CW-complex
\[
X = (C_0, C_1, C_2),
\]
where:
\begin{itemize}
  \item \(C_0\) is the set of 0-cells (vertices), representing agents, policies, and risk constraints.
  \item \(C_1\) is the set of 1-cells (edges), representing explicit bindings (technical specifications, contracts, SLAs, kill-switches).
  \item \(C_2\) is the set of 2-cells (faces), representing closed loops of accountability and verification.
\end{itemize}
We denote \(V = |C_0|\), \(E = |C_1|\), and \(F = |C_2|\).
\end{definition}

We allow a \emph{documented} lattice \(X_{\mathrm{doc}}\) and a \emph{runtime} lattice \(X_{\mathrm{run}}\) for the same organization, representing respectively the stated and realized structures.

\subsection{Euler Characteristic and Homology}

\begin{definition}[Euler characteristic]
For a finite CW-complex \(X\) with cells up to dimension \(2\), the Euler characteristic is
\[
\chi(X) = V - E + F.
\]
\end{definition}

\begin{definition}[Homology and Betti numbers]
Let \(H_k(X)\) be the \(k\)-th singular homology group of \(X\) (with coefficients in \(\mathbb{Z}\) or \(\mathbb{Q}\)). Its rank \(\beta_k = \mathrm{rank}\,H_k(X)\) is the \(k\)-th Betti number. Then
\[
\chi(X) = \sum_{k=0}^{\infty} (-1)^k \beta_k.
\]
In the 2-dimensional case:
\[
\chi(X) = \beta_0 - \beta_1 + \beta_2.
\]
\end{definition}

Interpretation in governance:
\begin{itemize}
  \item \(\beta_0\): number of connected components (silos).
  \item \(\beta_1\): number of non-contractible oversight loops (structural contradictions or unresolved cycles).
  \item \(\beta_2\): number of coverage voids (2D “holes” where oversight domains fail to meet).
\end{itemize}

\section{Gurau Degree for Governance Tensor Graphs}

\subsection{Definition}

We map a governance lattice onto a rank-\(d\) colored tensor graph as in Group Field Theory.

\begin{definition}[Decision vertices and internal faces]
Fix a dimension parameter \(d \ge 2\). Let \(V\) denote the number of \emph{decision vertices} (high-stakes decision events or commitment points). Let \(F_{\mathrm{int}}\) denote the number of \emph{internal faces} in the colored tensor graph corresponding to closed accountability loops internal to the governance manifold (no external boundary crossings).
\end{definition}

\begin{definition}[Governance Gurau degree]
The \emph{Gurau degree} of a governance tensor graph is
\[
G = d\,(V - 1) - (d - 1)\,F_{\mathrm{int}}.
\]
\end{definition}

Interpretation:
\begin{itemize}
  \item Negative or zero \(G\): amplitudes are well-behaved (manifold-like behavior; suppressed branched-polymer divergence).
  \item Positive \(G\): amplitudes diverge (branched-polymer regime; tree-like, fragile governance).
\end{itemize}

We will distinguish:
\[
G_{\mathrm{doc}} = d\,(V_{\mathrm{doc}} - 1) - (d-1)\,F_{\mathrm{int}}^{\mathrm{doc}},\quad
G_{\mathrm{run}} = d\,(V_{\mathrm{run}} - 1) - (d-1)\,F_{\mathrm{int}}^{\mathrm{run}}.
\]

\section{Structural Results}

\subsection{Monotonicity in Internal Faces}

\begin{lemma}[Monotonicity of \(G\) in \(F_{\mathrm{int}}\)]
Fix \(d > 1\) and \(V\). For two governance graphs with internal faces \(F_{\mathrm{int}}^{(1)},F_{\mathrm{int}}^{(2)}\) and corresponding degrees \(G^{(1)},G^{(2)}\), if
\[
F_{\mathrm{int}}^{(2)} > F_{\mathrm{int}}^{(1)},
\]
then
\[
G^{(2)} < G^{(1)}.
\]
\end{lemma}

\begin{proof}
By definition:
\[
G^{(i)} = d(V - 1) - (d-1)F_{\mathrm{int}}^{(i)}\quad (i=1,2).
\]
Then
\[
G^{(2)} - G^{(1)} = -(d-1)\big(F_{\mathrm{int}}^{(2)} - F_{\mathrm{int}}^{(1)}\big).
\]
Since \(d-1 > 0\) and \(F_{\mathrm{int}}^{(2)} - F_{\mathrm{int}}^{(1)} > 0\), we get \(G^{(2)} - G^{(1)} < 0\), i.e.\ \(G^{(2)} < G^{(1)}\).
\end{proof}

\begin{corollary}[Loop enrichment improves degree]
For fixed \(d,V\), increasing the number of internal faces strictly decreases the Gurau degree. Adding accountability loops (faces) makes the governance manifold more negative-\(G\), moving it away from the branched-polymer regime.
\end{corollary}

\subsection{Four-Vertex Necklace vs.\ Tree}

We specialize to \(d = 4\), as in the quantum spacetime analogy.

\begin{lemma}[Four-vertex tree]
If \(d=4\), \(V=4\), and \(F_{\mathrm{int}} \approx 0\) (a tree-like governance structure with essentially no closed loops), then
\[
G_{\mathrm{tree}} = 4(4-1) - 3\cdot 0 = 12.
\]
\end{lemma}

\begin{proof}
Direct substitution.
\end{proof}

\begin{lemma}[Four-vertex necklace]
If \(d=4\), \(V=4\), and \(F_{\mathrm{int}} = 9\) (multiple local loops and at least one global loop across the four decisions), then
\[
G_{\mathrm{necklace}} = 4(4-1) - 3\cdot 9 = 12 - 27 = -15.
\]
\end{lemma}

\begin{proof}
Again by direct substitution.
\end{proof}

\begin{corollary}[Topology and pathology]
With \(V=4\), the tree configuration yields \(G>0\) (branched-polymer regime), whereas the necklace configuration yields \(G\ll 0\) (manifold-like regime). For fixed \(V\), the difference is entirely due to the richer face structure of the necklace.
\end{corollary}

\subsection{Effect of Face Deletion on \(\chi\) and \(G\)}

\begin{lemma}[Euler characteristic under face deletion]
Let \(X\) be a 2-dimensional CW-complex with data \((V,E,F)\) and Euler characteristic \(\chi(X) = V - E + F\). If we obtain \(X'\) by removing a single 2-cell (face) without changing \(V\) or \(E\), so \(F' = F-1\), then
\[
\chi(X') = \chi(X) - 1.
\]
\end{lemma}

\begin{proof}
We have
\[
\chi(X') = V - E + (F-1) = (V - E + F) - 1 = \chi(X) - 1.
\]
\end{proof}

\begin{lemma}[Gurau degree deviation in terms of face loss]
Suppose \(V_{\mathrm{doc}} = V_{\mathrm{run}} = V\) and
\[
\Delta F_{\mathrm{int}} = F_{\mathrm{int}}^{\mathrm{doc}} - F_{\mathrm{int}}^{\mathrm{run}}.
\]
Then
\[
G_{\mathrm{run}} - G_{\mathrm{doc}} = (d - 1)\,\Delta F_{\mathrm{int}}.
\]
\end{lemma}

\begin{proof}
By definition:
\[
G_{\mathrm{doc}} = d(V - 1) - (d - 1)F_{\mathrm{int}}^{\mathrm{doc}},\quad
G_{\mathrm{run}} = d(V - 1) - (d - 1)F_{\mathrm{int}}^{\mathrm{run}}.
\]
Subtract:
\[
G_{\mathrm{run}} - G_{\mathrm{doc}}
= -(d-1)\big(F_{\mathrm{int}}^{\mathrm{run}} - F_{\mathrm{int}}^{\mathrm{doc}}\big)
= (d-1)\big(F_{\mathrm{int}}^{\mathrm{doc}} - F_{\mathrm{int}}^{\mathrm{run}}\big)
= (d-1)\,\Delta F_{\mathrm{int}}.
\]
\end{proof}

\begin{corollary}[Face loss increases \(G\)]
If faces are lost at runtime relative to documentation, i.e.\ \(F_{\mathrm{int}}^{\mathrm{run}} < F_{\mathrm{int}}^{\mathrm{doc}}\), then \(\Delta F_{\mathrm{int}} > 0\) and
\[
G_{\mathrm{run}} > G_{\mathrm{doc}}.
\]
Loss of operational loops moves the governance lattice toward more positive (worse) Gurau degree.
\end{corollary}

\section{Linear Operator Model and Norm Bounds}

\subsection{Topology Vector Space}

We encode each governance lattice as a vector in \(\mathbb{R}^4\).

\begin{definition}[Topology vector]
For a governance lattice \(X\), define
\[
\mathbf{t}(X) =
\begin{pmatrix}
V\\
E\\
F\\
G
\end{pmatrix}
\in \mathbb{R}^4,
\]
where \(V,E,F\) are as above, and \(G\) is the Gurau degree (with fixed \(d\)).
\end{definition}

Let \(\|\cdot\|_p\) denote any standard vector norm on \(\mathbb{R}^4\) (e.g.\ \(\ell_1,\ell_2,\ell_\infty\)).

We define:
\[
\mathbf{t}_{\mathrm{doc}} = \mathbf{t}(X_{\mathrm{doc}}),\quad
\mathbf{t}_{\mathrm{run}} = \mathbf{t}(X_{\mathrm{run}}),
\]
and the deviation vector
\[
\Delta \mathbf{t} = \mathbf{t}_{\mathrm{doc}} - \mathbf{t}_{\mathrm{run}}.
\]

\subsection{Deviation Operator}

We introduce a linear operator that approximates how documented topology is realized at runtime.

\begin{definition}[Deviation operator]
Let \(L:\mathbb{R}^4 \to \mathbb{R}^4\) be a linear operator and define the predicted runtime vector as
\[
\mathbf{t}_{\mathrm{pred}} = L\,\mathbf{t}_{\mathrm{doc}}.
\]
The residual is
\[
\mathbf{r} = \mathbf{t}_{\mathrm{run}} - \mathbf{t}_{\mathrm{pred}}.
\]
\end{definition}

A simple but instructive choice is a diagonal scaling matrix:
\[
L =
\begin{pmatrix}
\alpha_V & 0 & 0 & 0\\
0 & \alpha_E & 0 & 0\\
0 & 0 & \alpha_F & 0\\
0 & 0 & 0 & \alpha_G
\end{pmatrix},
\]
with \(0 \le \alpha_\bullet \le 1\) reflecting the fraction of documented structure realized at runtime (e.g.\ \(\alpha_F\) = fraction of documented loops that genuinely close).

\subsection{Operator Norm and Stability Bounds}

\begin{definition}[Operator norm]
Given norms \(\|\cdot\|_p\) on domain and codomain, the operator norm of \(L\) is
\[
\|L\|_{p \to p} = \sup_{\mathbf{x}\neq 0} \frac{\|L\mathbf{x}\|_p}{\|\mathbf{x}\|_p}.
\]
\end{definition}

\begin{lemma}[Norm of diagonal operator in \(\ell_2\)]
For the diagonal operator \(L\) above and the \(\ell_2\) norm,
\[
\|L\|_{2\to 2} = \max\{|\alpha_V|,|\alpha_E|,|\alpha_F|,|\alpha_G|\}.
\]
\end{lemma}

\begin{proof}
In the diagonal case, the eigenvalues of \(L\) are exactly the diagonal entries \(\alpha_V,\alpha_E,\alpha_F,\alpha_G\). The \(\ell_2\) operator norm equals the spectral radius for a normal operator, hence the maximum absolute eigenvalue.
\end{proof}

\begin{lemma}[Deviation bound in terms of \(\|L - I\|\)]
For any \(\mathbf{t}_{\mathrm{doc}}\),
\[
\|\mathbf{t}_{\mathrm{run}} - \mathbf{t}_{\mathrm{doc}}\|_p
\le
\|L - I\|_{p\to p}\,\|\mathbf{t}_{\mathrm{doc}}\|_p + \|\mathbf{r}\|_p,
\]
where \(I\) is the identity operator.
\end{lemma}

\begin{proof}
We write:
\[
\mathbf{t}_{\mathrm{run}} - \mathbf{t}_{\mathrm{doc}}
= \big(L\mathbf{t}_{\mathrm{doc}} - \mathbf{t}_{\mathrm{doc}}\big) + \mathbf{r}
= (L - I)\mathbf{t}_{\mathrm{doc}} + \mathbf{r}.
\]
Taking norms and using the triangle inequality:
\[
\|\mathbf{t}_{\mathrm{run}} - \mathbf{t}_{\mathrm{doc}}\|_p
\le
\|(L - I)\mathbf{t}_{\mathrm{doc}}\|_p + \|\mathbf{r}\|_p
\le
\|L - I\|_{p\to p}\,\|\mathbf{t}_{\mathrm{doc}}\|_p + \|\mathbf{r}\|_p.
\]
\end{proof}

\begin{corollary}[Relative deviation bound]
If \(\|L - I\|_{p\to p} \le \varepsilon < 1\), then
\[
\frac{\|\mathbf{t}_{\mathrm{run}} - \mathbf{t}_{\mathrm{doc}}\|_p}{\|\mathbf{t}_{\mathrm{doc}}\|_p}
\le
\varepsilon + \frac{\|\mathbf{r}\|_p}{\|\mathbf{t}_{\mathrm{doc}}\|_p}.
\]
Thus the topology vectors are close whenever both the operator deviation and the residual are small.
\end{corollary}

\subsection{Face-Realization Model and Gurau Drift}

We can identify a specific structure for \(L\) that captures loss of faces.

Assume:
\begin{itemize}
  \item \(V\) is unchanged from documented to runtime (\(\alpha_V = 1\)).
  \item Edges may partially realize (\(\alpha_E \le 1\)).
  \item Faces are realized at fraction \(\alpha_F\in[0,1]\):
  \[
  F_{\mathrm{run}} = \alpha_F F_{\mathrm{doc}}.
  \]
\end{itemize}

Then, with fixed \(d\), the difference in Gurau degree obeys exactly:
\[
G_{\mathrm{run}} - G_{\mathrm{doc}} = (d-1)\big(F_{\mathrm{int}}^{\mathrm{doc}} - F_{\mathrm{int}}^{\mathrm{run}}\big),
\]
as proven above. In the diagonal model, this corresponds to:
\[
\alpha_G = 1 \quad\text{and}\quad G\ \text{updated from the actual realized}\ F_{\mathrm{int}}.
\]

\begin{lemma}[Gurau drift as norm-bounded perturbation]
Consider an effective map \(T:\mathbb{R}^4 \to \mathbb{R}^4\) that takes \((V,E,F,G_{\mathrm{doc}})\) to \((V,E,F_{\mathrm{run}},G_{\mathrm{run}})\), with \(V,E\) unchanged and \(\Delta F = F_{\mathrm{doc}} - F_{\mathrm{run}}\). The change in the topology vector satisfies
\[
\big\|\mathbf{t}_{\mathrm{run}} - \mathbf{t}_{\mathrm{doc}}\big\|_p
\ge
c_p\,|\Delta F|,
\]
for some constant \(c_p>0\) depending on \(d\) and the choice of norm, reflecting that nonzero loss of faces necessarily yields a nonzero deviation vector and a Gurau drift magnitude
\[
|G_{\mathrm{run}} - G_{\mathrm{doc}}| = (d-1)|\Delta F|.
\]
\end{lemma}

\begin{proof}
The difference vector is
\[
\mathbf{t}_{\mathrm{run}} - \mathbf{t}_{\mathrm{doc}} =
\begin{pmatrix}
0\\
0\\
-\,\Delta F\\
(d-1)\,\Delta F
\end{pmatrix}.
\]
For any norm \(\|\cdot\|_p\), this vector has norm proportional to \(|\Delta F|\). For example, in \(\ell_2\),
\[
\|\mathbf{t}_{\mathrm{run}} - \mathbf{t}_{\mathrm{doc}}\|_2 =
|\Delta F|\sqrt{(-1)^2 + (d-1)^2} = |\Delta F|\sqrt{1 + (d-1)^2},
\]
so we may set \(c_2 = \sqrt{1 + (d-1)^2}\).
\end{proof}

\section{Worked Example: Two-Vertex Risk Loop Removal}

We make the preceding statements concrete for the primary illustrative example.

\subsection{Data}

Assume:
\begin{itemize}
  \item \(d = 4\).
  \item \(V = 2\) decision vertices.
  \item Documented faces: \(F_{\mathrm{int}}^{\mathrm{doc}} = 4\) (Reg–Legal, Technical–Controls, Risk–Capital, Ethics–Patient).
\end{itemize}

Documented Gurau degree:
\[
G_{\mathrm{doc}} = 4(2 - 1) - 3\cdot 4 = 4 - 12 = -8.
\]

\subsection{Runtime Loss}

Suppose at runtime the Risk–Capital loop fails to materialize:
\begin{itemize}
  \item \(F_{\mathrm{int}}^{\mathrm{run}} = 3\).
  \item \(\Delta F_{\mathrm{int}} = F_{\mathrm{int}}^{\mathrm{doc}} - F_{\mathrm{int}}^{\mathrm{run}} = 1.\)
\end{itemize}

Gurau drift:
\[
G_{\mathrm{run}} - G_{\mathrm{doc}} = (d-1)\Delta F_{\mathrm{int}} = 3\cdot 1 = 3,
\]
so
\[
G_{\mathrm{run}} = -8 + 3 = -5.
\]

Topology vector drift:
\[
\mathbf{t}_{\mathrm{doc}} =
\begin{pmatrix}
V\\E\\F_{\mathrm{doc}}\\G_{\mathrm{doc}}
\end{pmatrix}
=
\begin{pmatrix}
2\\E\\4\\-8
\end{pmatrix},\quad
\mathbf{t}_{\mathrm{run}} =
\begin{pmatrix}
2\\E\\3\\-5
\end{pmatrix},
\]
so
\[
\mathbf{t}_{\mathrm{run}} - \mathbf{t}_{\mathrm{doc}} =
\begin{pmatrix}
0\\0\\-1\\3
\end{pmatrix}.
\]

In \(\ell_2\):
\[
\|\mathbf{t}_{\mathrm{run}} - \mathbf{t}_{\mathrm{doc}}\|_2 = \sqrt{1^2 + 3^2} = \sqrt{10}.
\]

This quantifies the structural effect of deleting a single face: a fixed-magnitude jump in both Euler- and Gurau-related coordinates, and a nontrivial norm of the deviation vector even in this minimal case.

\section{Summary of Key Mathematical Facts}

The mathematical content of this appendix can be distilled into the following key facts:

\begin{enumerate}
  \item The Euler characteristic and Betti numbers of governance cell complexes detect holes and contradictions in the governance manifold.
  \item The Gurau degree, transplanted from colored tensor models, provides a scalar diagnostic of manifold-like vs.\ branched-polymer regimes for governance lattices.
  \item Increasing the number of internal faces at fixed decision-vertex count strictly decreases \(G\) (Lemma 1), while losing faces strictly increases \(G\) (Lemma 4).
  \item In a 4-vertex example, necklace and tree configurations with the same \(V\) can be sharply distinguished by their Gurau degrees (\(-15\) vs.\ \(12\)).
  \item The difference between documented and runtime governance lattices can be modeled as a linear-operator perturbation in \(\mathbb{R}^4\), with operator norms providing stability bounds for topological drift.
  \item Loss of a single accountability loop yields a computable, nonzero deviation in both Euler characteristic and Gurau degree, and a corresponding lower-bound on the norm of the topology drift vector.
\end{enumerate}

These results collectively provide a rigorous backbone for the Phase Mirror prescriptions articulated in the main text and protect them as explicit, formal prior art.

\nocite{*}
\bibliographystyle{unsrt}  
\bibliography{references}
\end{document}